\chapter{Impulse Control Problems}

\section*{Definition}
Impulse control problems involve a pattern of failing to resist urges, drive, or temptations to perform acts that are harmful to self or others. They are characterized by a sense of increasing tension before the act and pleasure, gratification, or relief during and after it.

\section*{Pathophysiology}
Impulse control problems reflect dysfunction in the fronto-striatal circuitry, particularly the orbitofrontal cortex, anterior cingulate cortex, and ventral striatum. Reduced prefrontal top-down inhibition combined with exaggerated ventral striatal reward signaling leads to a bias toward immediate gratification over long-term consequences.

\section*{Etiology}
\begin{itemize}
  \item \textbf{Psychiatric:} Intermittent explosive disorder, kleptomania, pyromania, trichotillomania, gambling disorder, conduct disorder, antisocial personality disorder, borderline personality disorder
  \item \textbf{Neurological:} Frontotemporal dementia, traumatic brain injury (orbitofrontal), Huntington disease, stroke
  \item \textbf{Medication-induced:} Dopamine agonists (pramipexole, ropinirole) used in Parkinson disease and restless legs syndrome; stimulants (methylphenidate, amphetamine)
  \item \textbf{Developmental:} Attention-deficit/hyperactivity disorder (ADHD), autism spectrum disorder
  \item \textbf{Psychosocial:} Childhood maltreatment, exposure to violence, poverty, poor parental supervision
\end{itemize}

\section*{Indications for Evaluation}
Seek care if:
\begin{itemize}
  \item Impaired control over impulses leads to harm (legal, financial, relational, physical)
  \item There is new onset of impulse control symptoms in the context of Parkinson disease treatment or head injury
  \item Intermittent explosive outbursts cause property damage or physical harm
\end{itemize}

\section*{Related Symptoms}
\begin{itemize}
  \item Anger outbursts, pathological gambling, compulsive shopping, kleptomania, skin picking, hair pulling, hypersexuality, substance use disorders, ADHD
\end{itemize}
\textbf{Note:} Impulse control disorders are underdiagnosed; screening is essential in patients taking dopamine agonists.

\section*{Self-Care / Home Management}
\begin{itemize}
  \item Remove access to triggers: limit credit cards and cash, block gambling websites, disable online shopping accounts
  \item Use the "urge surfing" technique: observe the urge without acting on it until it passes
  \item Create a list of alternative activities to engage in when the impulse arises
  \item Enlist family support for accountability
\end{itemize}

\section*{Diagnostic Approach}
\begin{itemize}
  \item Clinical interview using DSM-5 criteria for specific impulse control disorders
  \item Screener: Minnesota Impulse Disorders Interview (MIDI)
  \item Assess for comorbid ADHD, mood disorders, substance use, and personality disorders
  \item For Parkinson patients on dopamine agonists: Questionnaire for Impulsive-Compulsive Disorders in Parkinson's Disease (QUIP)
\end{itemize}

\section*{Treatment / Management Overview}
\begin{itemize}
  \item First-line psychotherapy: Cognitive-behavioral therapy (CBT) with relapse prevention
  \item For dopamine agonist-induced symptoms: reduce dose or switch to a different medication
  \item Pharmacotherapy: SSRIs (fluoxetine, citalopram) for trichotillomania and pathological gambling; naltrexone or nalmefene for gambling disorder
  \item Mood stabilizers (lithium, valproate) for intermittent explosive disorder
  \item Mutual help groups: Gamblers Anonymous, Debtors Anonymous
\end{itemize}

\clearpage
